\chapter*{Abstract}
\addcontentsline{toc}{chapter}{Abstract}

This textbook presents \emph{DAGonStar} (Direct Acyclic Graph On anything, also
written \emph{DAGon*}), a lightweight Python workflow engine designed to
orchestrate computational tasks across local machines, on-premise
high-performance computing (HPC) clusters, containers, cloud-based virtual
infrastructures, and, in its current evolution, artificial-intelligence (AI)
and web services. The material is the companion textbook of a 48-lecture course in which
theory and hands-on work are mixed in every lecture. It is aligned with
DAGonStar~3.0.0 (\url{https://github.com/DagOnStar/dagonstar}).

The book follows the four parts of the course. Part~I (Lectures~01--12)
establishes the foundations of scientific workflow orchestration: directed
acyclic graphs, the \texttt{workflow://} schema, configuration, scratch
folders, cycle detection and a first station pipeline. Part~II
(Lectures~13--24) covers execution, data and recovery: staging, garbage
collection, checkpointing, failures, native tasks, asynchronous runs,
meta-workflows, sweeps and ensembles. Part~III (Lectures~25--36) presents the
task models for AI--HPC integration: Web, LLM, FaaS, IoT and agentic tasks,
joined across the computing continuum, and an environmental case study.
Part~IV (Lectures~37--48) turns to infrastructure, data movement and FAIR
practice: containers, orchestrators, remote and cloud hosts, Slurm, BeeOND,
pre-staging, data movers, FAIR metadata, CWL export, operations and a
capstone. Appendices~A--D collect the API reference, developer practices,
configuration reference and checkpoint semantics. Appendix~E presents
DAGonFS, an in-memory MPI-distributed file system for workflow scratch data.
Appendices~F--H cover three additions around DAGonStar~3.0: ad-hoc volumes
shared by volume tasks and streams, meta-tasks that run workflows written for
other engines, and DAGonWeb, the web interface for designing and running
DAGonStar workflows.

Every course example is included verbatim from the \texttt{examples/}
folder of the course repository, together with its recorded output. The
script \texttt{tools/run\_examples.py} re-runs every example against the
installed DAGonStar and compares the outputs, so the book stays consistent
with the code.
